\section{Planets Queries I}
\label{sec:cses-planets-queries-i}

\problemstatement{Each of $n$ planets has exactly one teleporter to a
successor planet. Answer $q$ queries: starting from planet $x$ and
teleporting $k$ times (where $k$ can be up to $10^9$), which planet do you
reach?}

\begin{patternbox}
Following one edge per node many times is answered by \textbf{binary
lifting}: precompute where you land after $2^0,2^1,2^2,\ldots$ steps, then
compose the jumps for the set bits of $k$ to answer each query in
$O(\log k)$.
\end{patternbox}

\subsection*{Doubling the jump}

Let $\textit{up}[j][x]$ be the planet reached after $2^j$ teleports from
$x$. The base is $\textit{up}[0][x]=\textit{succ}(x)$, and doubling gives
$\textit{up}[j][x]=\textit{up}[j-1][\textit{up}[j-1][x]]$ -- jump $2^{j-1}$,
then $2^{j-1}$ more. For a query $(x,k)$, walk through the bits of $k$: for
each set bit $j$, replace $x$ with $\textit{up}[j][x]$. After processing all
bits, $x$ is the destination.

\begin{ideabox}[Any k Is a Sum of Powers of Two]
Because $k$ decomposes into its binary digits, a $k$-step jump is the
composition of the precomputed $2^j$-jumps for the bits set in $k$. Thus
$O(n\log k)$ preprocessing yields $O(\log k)$ per query -- the same doubling
trick as sparse tables and LCA.
\end{ideabox}

\subsection*{Algorithm}

\begin{enumerate}
  \item Read $\textit{succ}$; build $\textit{up}[0]$; fill
        $\textit{up}[j]$ by composing the previous level.
  \item For each $(x,k)$: for every set bit $j$ of $k$, advance $x=
        \textit{up}[j][x]$.
  \item Output the final $x$.
\end{enumerate}

\subsection*{Dry run}

\dryrun{$\textit{succ}=[2,3,1]$ (a 3-cycle); query $x=1,k=4$.}

\begin{itemize}
  \item $k=4=100_2$: jump $2^2=4$ steps. $\textit{up}[2][1]=$ four steps
        around the 3-cycle $=1\to2\to3\to1\to2$, landing on $2$.
  \item Answer $2$.
\end{itemize}

\subsection*{C++ implementation}

\begin{lstlisting}
#include <bits/stdc++.h>
using namespace std;
const int LOG = 30;                              // 2^30 > 1e9

int main() {
    int n, q;
    cin >> n >> q;
    vector<vector<int>> up(LOG, vector<int>(n + 1));
    for (int i = 1; i <= n; i++) cin >> up[0][i];  // successor
    for (int j = 1; j < LOG; j++)
        for (int i = 1; i <= n; i++)
            up[j][i] = up[j - 1][up[j - 1][i]];

    while (q--) {
        int x, k;
        cin >> x >> k;
        for (int j = 0; j < LOG; j++)
            if (k & (1 << j)) x = up[j][x];        // apply 2^j-jump for set bits
        cout << x << "\n";
    }
    return 0;
}
\end{lstlisting}

\begin{complexitybox}
\textbf{Time Complexity.} $O(n\log k)$ preprocessing, $O(\log k)$ per query.
\textbf{Space Complexity.} $O(n\log k)$ for the lifting table.
\end{complexitybox}

\begin{takeawaybox}
Binary lifting answers ``$k$ steps ahead'' in a functional graph in
logarithmic time by composing power-of-two jumps. The same table underlies
$k$-th ancestor queries and LCA, and powers the harder Planets Queries II
next.
\end{takeawaybox}
